\phantomsection
\section{Conclusion}
\label{sec:conclusion}

We have introduced TOBYQA, a time-augmented framework for derivative-free optimization under noise and temporal drift.

The analysis establishes solvability of the augmented KKT system under a full-column-rank condition on $X_{\mathrm{aug}}$. Because the composite kernel $K=A+\rho^{-1}I_m$ is positive definite, this condition is sufficient for a unique solution without the strict interpolation poisedness required by classical minimum-Frobenius-norm models (Proposition~\ref{prop:solvability}). The same recovery operator satisfies $\Pi X_{\mathrm{aug}}=I_{n+2}$. Consequently, under affine observation drift, the recovered spatial gradient is invariant to the drift magnitude, while under $C^2$ nonlinear drift its bias is bounded by $O(L_\mu\tau_{\max}^2)$ and therefore depends on the temporal span of the active sample set (Theorems~\ref{thm:exact-decoupling} and~\ref{thm:nonlinear-drift}).

The kernel also admits a statistical interpretation. Under a rotation-invariant Gaussian curvature prior, Powell's residual kernel $A_{ij}=\tfrac{1}{2}(\vct{d}_i^\top\vct{d}_j)^2$ is proportional to the covariance of the omitted quadratic forms, so the soft-constrained KKT estimate coincides with the generalized least-squares estimator under the resulting residual model (Proposition~\ref{prop:kernel-stat}). Under the fully-linear framework and the regularization confinement condition, the oracle complexity above the statistical noise floor is $O(L_g(f(\vct{x}_0)-f_{\mathrm{low}})\varepsilon^{-2})$, matching the lower bound for first-order nonconvex optimization (Theorem~\ref{thm:main-complexity}).

The experiments cover $65$ benchmark problems across three dimensions ($n \in \{6, 10, 20\}$) and seven evaluation channels. TOBYQA has the highest reported solve rates in these tests, with the largest differences occurring under additive drift. The ablations show lower solve rates when either drift compensation or the curvature prior kernel is removed. The evaluation-budget comparison also shows that the effect of fitting explicit quadratic curvature depends on both dimension and drift channel: it is less effective in low dimension and under nonlinear drift, but can be beneficial in higher dimension under the tested static and affine-drift settings. This pattern is not explained by the identifiability ratio of the quadratic coefficients alone.

Several questions remain open. The dimension-dependent behavior of the interpolated curvature variant calls for a criterion that predicts when curvature estimation is useful at the budget $m=2n+1$. The present framework assumes a static latent objective with a drifting observation channel; extending it to $f(\vct{x},t)$ would require a model for the motion of $\vct{x}_*(t)$ and corresponding tracking-error bounds. Polynomial time columns $[\vct{\theta},\vct{\theta}^2,\dots,\vct{\theta}^p]$ could be used to represent polynomial drift, but the conditioning of the enlarged system and its effect on recovery error remain to be analyzed. Extensions to bound constraints $\vct{l}\le\vct{x}\le\vct{u}$ and general linear constraints provide another direction for further work.
